\documentclass{article}
\usepackage{hyperref}
\usepackage{Style}

\nocite{*} % Comentar si quiero citar
%\addbibresource{bibliografia.bib} % Quitar el comentado si quiero usar bibliografia

\begin{document}

\begin{minipage}{2.5cm}
    \includegraphics[width=2cm]{imagen_puc.jpg}
\end{minipage}
\begin{minipage}{12cm}
    {\sc Pontificia Universidad Católica de Chile\\
    Facultad de Matemáticas\\
    Departamento de Matemática\\
    Profesor: Manuel Cabezas -- Ayudante: Benjamín Mateluna}
\end{minipage}
\vspace{2ex}

{\centerline{\bf Introducción a la Geometría - MAT1304}
\centerline{\bf Ayudantía 16}}
\centerline{\bf 05 de mayo de 2026}

\vspace{1cm}
\noindent\textbf{Problema 1.} Demuestre que para $z\in\C$, se tiene que
\begin{enumerate}
    \item $z\in\R$ si y solo si $z=\overline{z}$.
    \item $\abs{z}=1$ y $z\neq0$ entonces $\overline{z}=\frac{1}{z}$.
\end{enumerate}

\vspace{5mm}
\noindent\textbf{Problema 2.} Sea $z\in\C$ con parte imaginaria no nula tal que $z+\frac{1}{z}$ es 
un número real. Calcule $z\overline{z}$.

\vspace{5mm}
\noindent\textbf{Problema 3.} Sea $z=cis(\frac{2\pi}{5})$.
\begin{enumerate}
    \item Desmuestre que $z^{4}+z^{3}+z^{2}+z^{1}+1=0$.
    
    \item Pruebe que se cumple la siguiente expresión
    \begin{equation*}
        \left(z^{2}+\frac{1}{z^{2}}\right)+\left(z+\frac{1}{z}\right)+1=0
    \end{equation*}
    
    \item Verifique que $4cos^{2}(\frac{2\pi}{5})+2cos(\frac{2\pi}{5})-1=0$.
    
    \item Deduzca que
    \begin{equation*}
        cos\left(\frac{2\pi}{5}\right)=\frac{\sqrt{5}-1}{4}
    \end{equation*}
\end{enumerate}

\vspace{5mm}
\noindent\textbf{Problema 4.} Demuestre que todo polinomio en $\R[x]$ se puede escribir como 
multiplicación de polinomios lineales o cuadráticos con discriminante negativo.

\vspace{1cm}
\noindent\textbf{\underline{Ejercicios Propuestos:}}

\vspace{5mm}
\noindent\textbf{Extra 1.} Grafique en el plano complejo el conjunto
\begin{equation*}
    \mathcal{L}=\{z\in\C:Re(z)+Im(\overline{z})=Im(z)\}
\end{equation*}

\vspace{5mm}
\noindent\textbf{Extra 2.} Halle un polinomio $p(x)\in\R[x]$ de grado 5, con raíces $-1+i,
\sqrt{2}i$ y $1$ tal que $p(-1)=-6$.

%\printbibliography % Quitar el comentado si quiero usar bibliografia

\end{document}
